\chapter{2009年重庆卷（文）}

\section{单选题}

\begin{problem}
圆心在 \(y\) 轴上，半径为 \(1\)，且过点 \((1,2)\) 的圆的方程为（\quad）

\choices
    {\(x^{2} + (y - 2)^{2} = 1\)}
    {\(x^{2} + (y + 2)^{2} = 1\)}
    {\((x - 1)^{2} + (y - 3)^{2} = 1\)}
    {\(x^{2} + (y - 3)^{2} = 1\)}
\end{problem}

\begin{answer}
A.
\end{answer}

\begin{solution}
设圆心为 \(O(0,k)\)，由圆的半径为 \(1\) 且过点 \((1,2)\)，得
\[
1+(2-k)^2=1
\]
所以 \(k=2\)，圆的方程为 \(x^2+(y-2)^2=1\)，故选 A.
\end{solution}

\begin{problem}
命题“若一个数是负数，则它的平方是正数”的逆命题是（\quad）

\choices
    {“若一个数是负数，则它的平方不是正数”}
    {“若一个数的平方是正数，则它是负数”}
    {“若一个数不是负数，则它的平方不是正数”}
    {“若一个数的平方不是正数，则它不是负数”}
\end{problem}

\begin{answer}
B.
\end{answer}

\begin{solution}
原命题的条件是“一个数是负数”，结论是“它的平方是正数”. 交换条件和结论，得到逆命题“若一个数的平方是正数，则它是负数”，故选 B.
\end{solution}

\begin{problem}
\((x + 2)^{6}\) 的展开式中 \(x^{3}\) 的系数是（\quad）

\choices
    {\(20\)}
    {\(40\)}
    {\(80\)}
    {\(160\)}
\end{problem}

\begin{answer}
D.
\end{answer}

\begin{solution}
在二项式 \((x+2)^6\) 中，含有 \(x^3\) 的项是从 \(x^6\) 中取出三个 \(x\)，从其余三个因子中取出三个 \(2\) 所得到的项，其系数为
\[
\mathrm C_6^3 2^3=20\times 8=160
\]
故选 D.
\end{solution}

\begin{problem}
已知向量 \(\bs{a}=(1,1)\)，\(\bs{b}=(2,x)\). 若 \(\bs{a}+\bs{b}\) 与 \(4\bs{b}-2\bs{a}\) 平行，则实数 \(x\) 的值是（\quad）

\choices
    {\(-2\)}
    {\(0\)}
    {\(1\)}
    {\(2\)}
\end{problem}

\begin{answer}
D.
\end{answer}

\begin{solution}
由题意，
\(\bs{a}+\bs{b}=(3,x+1)\)，\(4\bs{b}-2\bs{a}=(6,4x-2)\).
两向量平行，所以其坐标行列式为零，即
\[
3(4x-2)-6(x+1)=0
\]
解得 \(6x-12=0\)，故 \(x=2\)，选 D.
\end{solution}

\begin{problem}
设 \(\{a_{n}\}\) 是公差不为 \(0\) 的等差数列，\(a_{1} = 2\) 且 \(a_{1}\)，\(a_{3}\)，\(a_{6}\) 成等比数列，则 \(\{a_{n}\}\) 的前 \(n\) 项和 \(S_{n} =\)（\quad）

\choices
    {\(\frac{n^2}{4} +\frac{7n}{4}\)}
    {\(\frac{n^2}{3} +\frac{5n}{3}\)}
    {\(\frac{n^2}{2} +\frac{3n}{4}\)}
    {\(n^2 + n\)}
\end{problem}

\begin{answer}
A.
\end{answer}

\begin{solution}
设等差数列的公差为 \(d\)，则
\(a_3=2+2d\)，\(a_6=2+5d\).
由 \(a_1\)，\(a_3\)，\(a_6\) 成等比数列，得
\[
(2+2d)^2=2(2+5d)
\]
即 \(4d^2-2d=0\). 因为 \(d\ne 0\)，所以 \(d=\frac12\). 于是
\[
a_n=2+\frac{n-1}{2}=\frac{n+3}{2}
\]
从而
\[
S_n=\frac n2(a_1+a_n)=\frac n2\left(2+\frac{n+3}{2}\right)=\frac{n^2}{4}+\frac{7n}{4}
\]
故选 A.
\end{solution}

\begin{problem}
下列关系式中正确的是（\quad）

\choices
    {\(\sin 11^{\circ} < \cos 10^{\circ} < \sin 168^{\circ}\)}
    {\(\sin 168^{\circ} < \sin 11^{\circ} < \cos 10^{\circ}\)}
    {\(\sin 11^{\circ} < \sin 168^{\circ} < \cos 10^{\circ}\)}
    {\(\sin 168^{\circ} < \cos 10^{\circ} < \sin 11^{\circ}\)}
\end{problem}

\begin{answer}
C.
\end{answer}

\begin{solution}
因为 \(168^{\circ}=180^{\circ}-12^{\circ}\)，所以
\[
\sin168^{\circ}=\sin12^{\circ}
\]
在 \((0^{\circ},90^{\circ})\) 内正弦函数单调递增，且
\(11^{\circ}<12^{\circ}<80^{\circ}\)，\(\cos10^{\circ}=\sin80^{\circ}\).
因此 \(\sin11^{\circ}<\sin168^{\circ}<\cos10^{\circ}\)，故选 C.
\end{solution}

\begin{problem}
已知 \(a > 0\)，\(b > 0\)，则 \(\frac{1}{a} + \frac{1}{b} + 2\sqrt{ab}\) 的最小值是（\quad）

\choices
    {\(2\)}
    {\(2\sqrt{2}\)}
    {\(4\)}
    {\(5\)}
\end{problem}

\begin{answer}
C.
\end{answer}

\begin{solution}
由基本不等式，
\[
\frac1a+\frac1b\geqslant \frac{2}{\sqrt{ab}}
\]
令 \(t=\sqrt{ab}>0\)，则原式满足
\[
\frac1a+\frac1b+2\sqrt{ab}
\geqslant \frac2t+2t\geqslant 4
\]
当且仅当 \(a=b\) 且 \(t=1\)，即 \(a=b=1\) 时取等号. 因此最小值为 \(4\)，选 C.
\end{solution}

\begin{problem}
\(12\) 个篮球队中有 \(3\) 个强队，将这 \(12\) 个队伍任意分成 \(3\) 个组（每组 \(4\) 个队），则 \(3\) 个强队恰好被分在同一组的概率为（\quad）

\choices
    {\(\frac{1}{55}\)}
    {\(\frac{3}{55}\)}
    {\(\frac{1}{4}\)}
    {\(\frac{1}{3}\)}
\end{problem}

\begin{answer}
B.
\end{answer}

\begin{solution}
固定包含一个强队的组. 其余 \(11\) 个队中，该组的三个空位等可能接纳其他队伍，因此另外两支强队都进入这个组的概率为
\[
\frac3{11}\times\frac2{10}=\frac3{55}
\]
故选 B.
\end{solution}

\begin{problem}
在正四棱柱 \(ABCD - A_{1}B_{1}C_{1}D_{1}\) 中，顶点 \(B_{1}\) 到对角线 \(BD_{1}\) 和到平面 \(A_{1}BCD_{1}\) 的距离分别为 \(h\) 和 \(d\)，则下列命题中正确的是（\quad）

\choices
    {若侧棱的长小于底面的边长，则 \(\frac{h}{d}\) 的取值范围为 \((0,1)\)}
    {若侧棱的长小于底面的边长，则 \(\frac{h}{d}\) 的取值范围为 \(\left(\frac{\sqrt{2}}{2}, \frac{2\sqrt{3}}{3}\right)\)}
    {若侧棱的长大于底面的边长，则 \(\frac{h}{d}\) 的取值范围为 \(\left(\frac{2\sqrt{3}}{3}, \sqrt{2}\right)\)}
    {若侧棱的长大于底面的边长，则 \(\frac{h}{d}\) 的取值范围为 \(\left(\frac{2\sqrt{3}}{3}, +\infty\right)\)}
\end{problem}

\begin{answer}
C.
\end{answer}

\begin{solution}
设正四棱柱的底面边长为 \(a\)，侧棱长为 \(l\). 建立空间直角坐标系，令
\(A(0,0,0)\)，\(B(a,0,0)\)，\(C(a,a,0)\)，\(D(0,a,0)\)
\(A_1(0,0,l)\)，\(B_1(a,0,l)\)，\(D_1(0,a,l)\).
直线 \(BD_1\) 的方向向量为 \((-a,a,l)\)，而 \(\overrightarrow{BB_1}=(0,0,l)\)，所以点 \(B_1\) 到直线 \(BD_1\) 的距离为
\[
h=\frac{\left|\overrightarrow{BB_1}\times\overrightarrow{BD_1}\right|}{|\overrightarrow{BD_1}|}
=\frac{al\sqrt2}{\sqrt{2a^2+l^2}}
\]
平面 \(A_1BCD_1\) 的一个法向量可取 \((l,0,a)\)，其方程为
\[
lx+az=al
\]
因此
\[
d=\frac{|la+al-al|}{\sqrt{l^2+a^2}}=\frac{al}{\sqrt{l^2+a^2}}
\]
于是
\[
\frac hd=\frac{\sqrt2\sqrt{l^2+a^2}}{\sqrt{l^2+2a^2}}
\]
令 \(t=l/a>0\)，则
\[
\left(\frac hd\right)^2=\frac{2(t^2+1)}{t^2+2}
\]
并且
\[
\frac{\mathrm d}{\mathrm dt}\left(\frac{2(t^2+1)}{t^2+2}\right)
=\frac{4t}{(t^2+2)^2}>0
\]
所以该比值随 \(t\) 增大而增大. 当 \(l>a\)，即 \(t>1\) 时，
\[
\frac{2\sqrt3}{3}<\frac hd<\sqrt2
\]
故选 C.
\end{solution}

\begin{problem}
把函数 \( f(x) = x^3 - 3x \) 的图象 \(C_1\) 向右平移 \( u \) 个单位长度，再向下平移 \( v \) 个单位长度后得到图象 \(C_2\). 若对任意的 \( u > 0 \)，曲线 \(C_1\) 与 \(C_2\) 至多只有一个交点，则 \( v \) 的最小值为（\quad）

\choices
    {\(2\)}
    {\(4\)}
    {\(6\)}
    {\(8\)}
\end{problem}

\begin{answer}
B.
\end{answer}

\begin{solution}
向右平移 \(u\) 个单位、向下平移 \(v\) 个单位后，曲线 \(C_2\) 的方程为
\[
y=f(x-u)-v
\]
交点的横坐标满足
\[
f(x)=f(x-u)-v
\]
即
\[
3ux^2-3u^2x+u^3-3u+v=0
\]
因为 \(u>0\)，要使任意 \(u>0\) 时至多有一个交点，必须且只需该关于 \(x\) 的二次方程判别式不大于零. 其判别式为
\[
\Delta=9u^4-12u(u^3-3u+v)=-3u\left(u^3-12u+4v\right)
\]
所以需对一切 \(u>0\) 有
\[
u^3-12u+4v\geqslant0
\]
即
\[
v\geqslant 3u-\frac{u^3}{4}
\]
函数 \(3u-\frac{u^3}{4}\) 在 \(u>0\) 上的导数为 \(3-\frac{3u^2}{4}\)，故其最大值在 \(u=2\) 处取得，最大值为 \(4\). 因此 \(v\) 的最小值为 \(4\)，选 B.
\end{solution}

\section{填空题}

\begin{problem}
\(U = \{n \mid n\) 是小于 \(9\) 的正整数\(\}\)，\(A = \{n \in U \mid n\) 是奇数\(\}\)，\(B = \{n \in U \mid n\) 是 \(3\) 的倍数\(\}\)，则 \(\mathrm{C}_U(A \cup B)=\)\fillinblank{}.
\end{problem}

\begin{answer}
\(\{2,4,8\}\).
\end{answer}

\begin{solution}
由题意，\(U=\{1,2,3,4,5,6,7,8\}\)，\(A=\{1,3,5,7\}\)，\(B=\{3,6\}\). 所以
\[
A\cup B=\{1,3,5,6,7\}
\]
其在 \(U\) 中的补集为
\[
\mathrm{C}_U(A\cup B)=\{2,4,8\}
\]
\end{solution}

\begin{problem}
记 \( f(x) = \log_{3}(x + 1) \) 的反函数为 \( y = f^{-1}(x) \)，则方程 \( f^{-1}(x) = 8 \) 的解 \( x = \)\fillinblank{}.
\end{problem}

\begin{answer}
\(2\).
\end{answer}

\begin{solution}
由 \(y=\log_3(x+1)\)，交换 \(x\)，\(y\) 得
\(x=\log_3(y+1)\)，\(y=3^x-1\).
因此 \(f^{-1}(x)=3^x-1\). 代入方程得
\[
3^x-1=8
\]
即 \(3^x=9=3^2\)，所以 \(x=2\).
\end{solution}

\begin{problem}
\(5\) 个人站成一排，其中甲，乙两人不相邻的排法有\fillinblank{}种. （用数字作答）
\end{problem}

\begin{answer}
\(72\).
\end{answer}

\begin{solution}
五个人的全排列共有 \(5!=120\) 种. 把甲、乙看作一个整体时，甲乙或乙甲有 \(2\) 种内部顺序，该整体与另外三人共有 \(4!\) 种排列，所以甲、乙相邻的排法有
\[
2\times4!=48
\]
种. 因此甲、乙不相邻的排法有
\[
5!-2\times4!=120-48=72
\]
种.
\end{solution}

\begin{problem}
从一堆苹果中任取 \(5\) 只，称得它们的质量如下（单位：\(\text{克}\)）
\begin{center}
\begin{tabular}{ccccc}
\(125\) & \(124\) & \(121\) & \(123\) & \(127\)
\end{tabular}
\end{center}
则该样本标准差 \(s =\)\fillinblank{}\(\text{克}\). （用数字作答）
\end{problem}

\begin{answer}
\(2\).
\end{answer}

\begin{solution}
五个质量数据的平均数为
\[
\bar x=\frac{125+124+121+123+127}{5}=124
\]
各数据与平均数的差分别为 \(1\)，\(0\)，\(-3\)，\(-1\)，\(3\)，所以
\[
\sum_{i=1}^{5}(x_i-\bar x)^2=1^2+0^2+(-3)^2+(-1)^2+3^2=20
\]
按样本标准差的定义
\[
s=\sqrt{\frac15\sum_{i=1}^{5}(x_i-\bar x)^2}=\sqrt{\frac{20}{5}}=2
\]
\end{solution}

\begin{problem}
已知椭圆 \(\frac{x^2}{a^2} + \frac{y^2}{b^2} = 1\)（\(a > b > 0\)）的左，右焦点分别为 \(F_1(-c,0)\)，\(F_2(c,0)\).
若椭圆上存在一点 \(P\) 使 \(\frac{a}{\sin \angle PF_1F_2} = \frac{c}{\sin \angle PF_2F_1}\)，则该椭圆的离心率的取值范围为\fillinblank{}.
\end{problem}

\begin{answer}
\(\left(\sqrt{2}-1,1\right)\).
\end{answer}

\begin{solution}
设椭圆的离心率为 \(e=\frac ca\)，则 \(0<e<1\). 在三角形 \(PF_1F_2\) 中，由正弦定理，题设条件给出
\[
\frac{PF_2}{PF_1}=\frac{\sin\angle PF_1F_2}{\sin\angle PF_2F_1}=\frac ac=\frac1e
\]
即
\[
\frac{PF_1}{PF_2}=e
\]
设 \(P=(x,y)\). 椭圆的焦半径公式为
\(PF_1=a+ex\)，\(PF_2=a-ex\).
因此
\[
\frac{a+ex}{a-ex}=e
\]
从而
\[
\frac xa=\frac{e-1}{e(1+e)}
\]
要使椭圆上存在这样的点，必须有 \(-1\leqslant x/a\leqslant1\). 由于 \(0<e<1\)，右端为负数，故只需
\[
\frac{1-e}{e(1+e)}\leqslant1
\]
整理得
\[
e^2+2e-1\geqslant0
\]
所以 \(e\geqslant\sqrt2-1\). 当 \(e=\sqrt2-1\) 时，\(x/a=-1\)，点 \(P\) 为椭圆左顶点，三点 \(P\)，\(F_1\)，\(F_2\) 共线，题设中的两个正弦均为零，原等式无意义，故应排除该端点. 当 \(e>\sqrt2-1\) 时，\(-1<x/a<0\)，由椭圆方程可取
\[
y=\pm b\sqrt{1-\frac{x^2}{a^2}}
\]
从而确有椭圆上的点 \(P\) 满足条件. 于是离心率的取值范围为
\[
\left(\sqrt2-1,1\right)
\]
\end{solution}

\section{解答题}

\begin{problem}
设函数 \( f(x) = (\sin \omega x + \cos \omega x)^2 + 2\cos^2 \omega x \)（\(\omega > 0\)）的最小正周期为 \( \frac{2\pi}{3} \).
\begin{enumerate}
\item 求 \( \omega \) 的值.
\item 若函数 \( y = g(x) \) 的图象是由 \( y = f(x) \) 的图象向右平移 \( \frac{\pi}{2} \) 个单位长度得到，求 \( y = g(x) \) 的单调增区间.
\end{enumerate}
\end{problem}

\begin{solution}
\begin{enumerate}
\item 由
\[
f(x)=(\sin\omega x+\cos\omega x)^2+2\cos^2\omega x
=2+\sin2\omega x+\cos2\omega x
=2+\sqrt2\sin\left(2\omega x+\frac\pi4\right)
\]
可知 \(f(x)\) 的最小正周期为 \(\frac{2\pi}{2\omega}=\frac\pi\omega\). 由题设
\[
\frac\pi\omega=\frac{2\pi}{3}
\]
所以 \(\omega=\frac32\).

\item 图象向右平移 \(\frac\pi2\) 个单位后，
\[
g(x)=f\left(x-\frac\pi2\right)=2+\cos(3x)-\sin(3x)
\]
因此
\[
g'(x)=-3\sin(3x)-3\cos(3x)
\]
\(g(x)\) 单调递增当且仅当
\[
\sin(3x)+\cos(3x)<0
\]
即
\[
\sqrt2\sin\left(3x+\frac\pi4\right)<0
\]
所以
\(3x+\frac\pi4\in\left(\pi+2k\pi,2\pi+2k\pi\right)\)，\(k\in\Z\)
从而单调增区间为
\(\left(\frac\pi4+\frac{2k\pi}{3},\frac{7\pi}{12}+\frac{2k\pi}{3}\right)\)，\(k\in\Z\).
\end{enumerate}
\end{solution}

\begin{problem}
某单位为绿化环境，移栽了甲，乙两种大树各 \(2\) 株. 设甲，乙两种大树移栽的成活率分别为 \(\frac{5}{6}\) 和 \(\frac{4}{5}\)，且各株大树是否成活互不影响. 求移栽的 \(4\) 株大树中：
\begin{enumerate}
\item 至少有 \(1\) 株成活的概率；
\item 两种大树各成活 \(1\) 株的概率.
\end{enumerate}
\end{problem}

\begin{solution}
\begin{enumerate}
\item 四株大树均未成活的概率为
\[
\left(1-\frac56\right)^2\left(1-\frac45\right)^2
=\left(\frac16\right)^2\left(\frac15\right)^2=\frac1{900}
\]
所以至少有一株成活的概率为
\[
1-\frac1{900}=\frac{899}{900}
\]

\item 甲种大树恰好成活一株的概率为
\[
\mathrm C_2^1\frac56\frac16=\frac5{18}
\]
乙种大树恰好成活一株的概率为
\[
\mathrm C_2^1\frac45\frac15=\frac8{25}
\]
由各株大树是否成活互不影响，两种大树各成活一株的概率为
\[
\frac5{18}\times\frac8{25}=\frac4{45}
\]
\end{enumerate}
\end{solution}

\begin{problem}
如图，在五面体 \(ABCDEF\) 中，\(AB \parallel DC\)，\(\angle BAD = \frac{\pi}{2}\)，\(CD = AD = 2\)，四边形 \(ABFE\) 为平行四边形，\(FA \perp\) 平面 \(ABCD\)，\(FC = 3\)，\(ED = \sqrt{7}\). 求：
\begin{enumerate}
\item 直线 \(AB\) 到平面 \(EFCD\) 的距离；
\item 二面角 \(F - AD - E\) 的平面角的正切值.
\end{enumerate}

\bitmapfigure[max width=0.56\linewidth]{img/2009/chongqing_liberal/q18_fig1.png}
\end{problem}

\begin{solution}
建立空间直角坐标系，使底面 \(ABCD\) 在 \(xOy\) 平面内，令
\(A(0,0,0)\)，\(B(p,0,0)\)，\(D(0,2,0)\)，\(C(2,2,0)\)
其中 \(p=AB>0\). 设 \(FA=q>0\)，由 \(FA\perp\) 平面 \(ABCD\) 及平行四边形 \(ABFE\)，可取
\(F(0,0,q)\)，\(E(-p,0,q)\).
由 \(FC=3\)，得
\[
FC^2=2^2+2^2+q^2=9
\]
所以 \(q=1\). 又由 \(ED=\sqrt7\)，得
\[
ED^2=p^2+2^2+q^2=7
\]
从而 \(p^2=2\)，即 \(p=\sqrt2\).

\begin{enumerate}
\item 平面 \(EFCD\) 经过 \(F(0,0,1)\)、\(D(0,2,0)\)，且平行于 \(x\) 轴，其方程为
\[
y+2z=2
\]
直线 \(AB\) 的方向为 \(x\) 轴方向，故 \(AB\parallel\) 平面 \(EFCD\). 于是所求距离等于点 \(A\) 到该平面的距离：
\[
d=\frac{|0+2\times0-2|}{\sqrt{1^2+2^2}}=\frac2{\sqrt5}
\]

\item 因为 \(FA\perp\) 平面 \(ABCD\)，所以 \(FA\perp AD\). 又 \(AD\perp AB\)，且 \(AB\parallel FE\)，故 \(AD\perp FE\). 由平行四边形 \(ABFE\)，有
\[
\overrightarrow{AE}=\overrightarrow{BF}=\overrightarrow{BA}+\overrightarrow{AF}
\]
而 \(AD\perp BA\) 且 \(AD\perp AF\)，所以 \(AD\perp AE\). 从而 \(AF\) 和 \(AE\) 分别是两个半平面内过 \(A\) 且垂直于棱 \(AD\) 的直线，所以 \(\angle FAE\) 是二面角 \(F-AD-E\) 的平面角.

又因为 \(AF\perp\) 平面 \(ABCD\)，而 \(FE\parallel AB\subset\) 平面 \(ABCD\)，所以 \(AF\perp FE\). 在直角三角形 \(AFE\) 中，
\[
\tan\angle FAE=\frac{FE}{AF}=\frac{AB}{FA}=\frac{\sqrt2}{1}=\sqrt2
\]
\end{enumerate}
\end{solution}

\begin{problem}
已知 \(f(x)=x^2+bx+c\) 为偶函数，曲线 \(y=f(x)\) 过点 \((2,5)\)，\(g(x)=(x+a)f(x)\).
\begin{enumerate}
\item 求曲线 \(y=g(x)\) 斜率为 \(0\) 的切线，求实数 \(a\) 的取值范围；
\item 若当 \(x=-1\) 时函数 \(y=g(x)\) 取得极值，确定 \(y=g(x)\) 的单调区间.
\end{enumerate}
\end{problem}

\begin{solution}
\begin{enumerate}
\item 因为 \(f(x)=x^2+bx+c\) 是偶函数，所以 \(b=0\). 又曲线过点 \((2,5)\)，故
\(4+c=5\)，\(c=1\).
于是 \(f(x)=x^2+1\)，从而
\[
g(x)=(x+a)(x^2+1)=x^3+ax^2+x+a
\]
曲线的切线斜率为 \(g'(x)=3x^2+2ax+1\). 存在斜率为 \(0\) 的切线，等价于方程
\[
3x^2+2ax+1=0
\]
有实数解，即
\[
\Delta=(2a)^2-4\times3\times1=4(a^2-3)\geqslant0
\]
所以 \(a\) 的取值范围为
\[
(-\infty,-\sqrt3]\cup[\sqrt3,+\infty)
\]

\item 当 \(x=-1\) 时函数取得极值，必须有 \(g'(-1)=0\). 因此
\(3-2a+1=0\)，\(a=2\).
此时
\[
g'(x)=3x^2+4x+1=(3x+1)(x+1)
\]
结合两个零点 \(-1\) 和 \(-\frac13\)，得 \(g'(x)>0\) 在
\[
(-\infty,-1)\cup\left(-\frac13,+\infty\right)
\]
上成立，\(g'(x)<0\) 在
\[
\left(-1,-\frac13\right)
\]
上成立. 因此 \(g(x)\) 的单调增区间为
\[
(-\infty,-1)\cup\left(-\frac13,+\infty\right)
\]
单调减区间为
\[
\left(-1,-\frac13\right)
\]
\end{enumerate}
\end{solution}

\begin{problem}
已知以原点 \(O\) 为中心的双曲线的一条准线方程为 \(x = \frac{\sqrt{5}}{5}\)，离心率 \(e = \sqrt{5}\).
\begin{enumerate}
\item 求该双曲线的方程.
\item 如图，点 \(A\) 的坐标为 \((- \sqrt{5}, 0)\)，\(B\) 是圆 \(x^{2} + (y - \sqrt{5})^{2} = 1\) 上的点，点 \(M\) 在双曲线右支上，求 \(|MA| + |MB|\) 的最小值，并求此时 \(M\) 点的坐标.
\end{enumerate}

\bitmapfigure[max width=0.42\linewidth]{img/2009/chongqing_liberal/q20_fig1.png}
\end{problem}

\begin{solution}
\begin{enumerate}
\item 设双曲线方程为
\[
\frac{x^2}{a^2}-\frac{y^2}{b^2}=1
\]
其焦半距为 \(c\). 由于准线方程为 \(x=\frac ae=\frac{\sqrt5}{5}\)，且 \(e=\sqrt5\)，所以
\[
a=e\cdot\frac{\sqrt5}{5}=1
\]
于是 \(c=ae=\sqrt5\)，并且
\[
b^2=c^2-a^2=5-1=4
\]
故双曲线方程为
\[
x^2-\frac{y^2}{4}=1
\]

\item 记双曲线右焦点为 \(F(\sqrt5,0)\)，圆心为 \(O_1(0,\sqrt5)\). 对双曲线右支上的点 \(M\)，有 \(x_M\geqslant1\). 若 \(x_M=1\)，由双曲线方程得 \(y_M=0\)，此时 \(MO_1=\sqrt6>1\)；若 \(x_M>1\)，则
\[
MO_1=\sqrt{x_M^2+(y_M-\sqrt5)^2}\geqslant x_M>1
\]
因此 \(M\) 在圆外. 由双曲线的定义，
\[
MA-MF=2a=2
\]
又因为圆的半径为 \(1\)，所以
\[
MB\geqslant MO_1-1
\]
等号在 \(B\) 是从 \(O_1\) 指向 \(M\) 的射线与圆的交点时取得. 因此
\[
MA+MB\geqslant MF+MO_1+1\geqslant FO_1+1=\sqrt{(\sqrt5)^2+(\sqrt5)^2}+1=\sqrt{10}+1
\]

下面说明等号可以取得. 取 \(M\) 在线段 \(FO_1\) 上，则其坐标满足 \(x+y=\sqrt5\). 将 \(y=\sqrt5-x\) 代入双曲线方程，得
\[
x^2-\frac{(\sqrt5-x)^2}{4}=1
\]
即
\[
3x^2+2\sqrt5x-9=0
\]
解得
\[
x=\frac{4\sqrt2-\sqrt5}{3}\quad\text{或}\quad x=-\frac{4\sqrt2+\sqrt5}{3}
\]
取在线段 \(FO_1\) 上的交点，得到
\[
M\left(\frac{4\sqrt2-\sqrt5}{3},\frac{4(\sqrt5-\sqrt2)}{3}\right)
\]
此时 \(M\) 在双曲线右支上，且 \(M\) 位于 \(F\)，\(O_1\) 之间；再取圆上满足 \(O_1B\) 与 \(O_1M\) 同向且 \(O_1B=1\) 的点，即可同时取得上述两个三角不等式的等号. 因此
\[
|MA|+|MB|\text{ 的最小值为 }\sqrt{10}+1
\]
\end{enumerate}
\end{solution}

\begin{problem}
已知 \(a_1 = 1\)，\(a_2 = 4\)，\(a_{n+2} = 4a_{n+1} + a_n\)，\(b_n = \frac{a_{n+1}}{a_n}\)，\(n \in \mathbf{N}^*\).
\begin{enumerate}
\item 求 \(b_{1}\)，\(b_{2}\)，\(b_{3}\) 的值；
\item 设 \(c_{n} = b_{n}b_{n + 1}\)，\(S_{n}\) 为数列 \(\{c_n\}\) 的前 \(n\) 项和. 求证： \(S_{n} \geqslant 17n\)；
\item 求证： \( |b_{2n} - b_n| < \frac{1}{64} \cdot \frac{1}{17^{n-2}} \).
\end{enumerate}
\end{problem}

\begin{solution}
\begin{enumerate}
\item 由递推关系，
\[
a_3=4a_2+a_1=4\times4+1=17
\]
\[
a_4=4a_3+a_2=4\times17+4=72
\]
所以
\(b_1=\frac{a_2}{a_1}=4\)，\(b_2=\frac{a_3}{a_2}=\frac{17}{4}\)，\(b_3=\frac{a_4}{a_3}=\frac{72}{17}\).

\item 由定义，
\[
c_n=b_nb_{n+1}=\frac{a_{n+1}}{a_n}\cdot\frac{a_{n+2}}{a_{n+1}}=\frac{a_{n+2}}{a_n}=4\frac{a_{n+1}}{a_n}+1=4b_n+1
\]
由 \(a_1=1>0\)、\(a_2=4>0\) 及递推式可归纳得 \(a_n>0\) 对一切 \(n\in\mathbf N^*\) 成立，且
\[
b_{n+1}=\frac{a_{n+2}}{a_{n+1}}=4+\frac{a_n}{a_{n+1}}>4
\]
结合 \(b_1=4\)，可知 \(b_n\geqslant4\)，于是 \(c_n=4b_n+1\geqslant17\). 因此
\[
S_n=\sum_{k=1}^{n}c_k\geqslant\sum_{k=1}^{n}17=17n
\]

\item 令
\[
D_k=a_{k+2}a_k-a_{k+1}^2
\]
由递推关系 \(a_k=a_{k+2}-4a_{k+1}\)，有
\[
D_{k+1}=a_{k+3}a_{k+1}-a_{k+2}^2
=(4a_{k+2}+a_{k+1})a_{k+1}-a_{k+2}^2
=-D_k
\]
又 \(D_1=a_3a_1-a_2^2=17-16=1\)，所以 \(|D_k|=1\). 从而
\[
|b_{k+1}-b_k|=\left|\frac{a_{k+2}}{a_{k+1}}-\frac{a_{k+1}}{a_k}\right|
=\frac{|D_k|}{a_{k+1}a_k}=\frac1{a_{k+1}a_k}
\]
记 \(P_k=a_{k+1}a_k\)，则
\[
\frac{P_{k+1}}{P_k}=\frac{a_{k+2}}{a_k}=b_kb_{k+1}=c_k\geqslant17
\]
因为 \(P_1=a_2a_1=4\)，所以
\[
P_k\geqslant4\cdot17^{k-1}
\]
于是
\[
|b_{2n}-b_n|\leqslant\sum_{k=n}^{2n-1}|b_{k+1}-b_k|=\sum_{k=n}^{2n-1}\frac1{P_k}
\]
又因为 \(P_k\geqslant4\cdot17^{k-1}\)，所以
\[
|b_{2n}-b_n|\leqslant\frac14\sum_{k=n}^{2n-1}\frac1{17^{k-1}}<\frac14\sum_{k=n}^{\infty}\frac1{17^{k-1}}=\frac1{64}\cdot\frac1{17^{n-2}}
\]
这就证明了所要结论.
\end{enumerate}
\end{solution}
